% Общая преамбула презентаций (beamer/metropolis).
% Теоремные окружения beamer определяет сам — здесь только их русские имена
% и те же макросы обозначений, что в конспектах.
\usepackage{amsmath,amssymb,mathtools}

\setbeamertemplate{theorems}[numbered]
\uselanguage{russian}
\languagepath{russian}
\deftranslation[to=russian]{Theorem}{Теорема}
\deftranslation[to=russian]{Lemma}{Лемма}
\deftranslation[to=russian]{Corollary}{Следствие}
\deftranslation[to=russian]{Definition}{Определение}
\deftranslation[to=russian]{Definitions}{Определения}
\deftranslation[to=russian]{Example}{Пример}
\deftranslation[to=russian]{Examples}{Примеры}
\deftranslation[to=russian]{Proof}{Доказательство}
\deftranslation[to=russian]{Problem}{Задача}
\deftranslation[to=russian]{Remark}{Замечание}

% beamer определяет theorem/lemma/corollary/definition/example/problem/proof сам,
% но не remark и не notation — добавляем, иначе сборка падает на \begin{remark}.
\theoremstyle{remark}
\newtheorem{remark}{Замечание}
\newtheorem{notation}{Обозначения}

\newcommand{\R}{\mathbb{R}}
\newcommand{\C}{\mathbb{C}}
\newcommand{\N}{\mathbb{N}}
\newcommand{\Z}{\mathbb{Z}}
\newcommand{\E}{\mathsf{E}}
\newcommand{\Var}{\mathsf{D}}
\newcommand{\Prob}{\mathsf{P}}
\newcommand{\T}{^{\mathsf{T}}}
\newcommand{\norm}[1]{\left\lVert #1 \right\rVert}
\newcommand{\abs}[1]{\left\lvert #1 \right\rvert}
\newcommand{\scal}[2]{\left\langle #1, #2 \right\rangle}
\DeclareMathOperator*{\argmin}{arg\,min}
\DeclareMathOperator*{\argmax}{arg\,max}
\DeclareMathOperator{\sgn}{sign}
\DeclareMathOperator{\rank}{rank}
\DeclareMathOperator{\tr}{tr}
\DeclareMathOperator{\diag}{diag}
\DeclareMathOperator{\supp}{supp}

% Многочлены Чебышёва: T_n — тригонометрическая нормировка cos(n arccos x),
% \Tm_n — приведённая (старший коэффициент 1). Нужны в вопросах 01, 11, 13, 19.
\newcommand{\Tm}{\widetilde{T}}
